\documentclass[12pt,a4paper]{article}

% ============================================================
% Packages
% ============================================================
\usepackage[utf8]{inputenc}
\usepackage[T1]{fontenc}
\usepackage{amsmath,amssymb,amsthm}
\usepackage{mathtools}
\usepackage{hyperref}
\usepackage[margin=1in]{geometry}
\usepackage{enumitem}
\usepackage{booktabs}
\usepackage{microtype}
\usepackage{xcolor}

% ============================================================
% Theorem environments
% ============================================================
\newtheorem{theorem}{Theorem}[section]
\newtheorem{lemma}[theorem]{Lemma}
\newtheorem{proposition}[theorem]{Proposition}
\newtheorem{corollary}[theorem]{Corollary}
\theoremstyle{definition}
\newtheorem{definition}[theorem]{Definition}
\newtheorem{axiom_def}[theorem]{Axiom}
\newtheorem{remark}[theorem]{Remark}

% ============================================================
% Custom commands
% ============================================================
\newcommand{\lean}[1]{\texttt{#1}}
\newcommand{\Hfunc}{\mathcal{H}}
\newcommand{\R}{\mathbb{R}}
\newcommand{\N}{\mathbb{N}}
\newcommand{\NNR}{\mathbb{R}_{\geq 0}}

% ============================================================
% Hyperref configuration
% ============================================================
\hypersetup{
  colorlinks=true,
  linkcolor=blue!70!black,
  citecolor=green!50!black,
  urlcolor=blue!70!black,
  pdfauthor={Essam Abadir},
  pdftitle={The Man Who Understood Entropy},
}

% ============================================================
\begin{document}

% ============================================================
% Title
% ============================================================
\title{%
  \textbf{The Man Who Understood Entropy:}\\[6pt]
  \large Correcting the Scientific Record on the Biggest\\
  Open Question in Science%
}
\author{Essam Abadir}
\date{March 2026}
\maketitle

% ============================================================
% Dedication
% ============================================================
\begin{center}
  \textit{In memory of Gian-Carlo Rota (1932--1999), who understood entropy.}
\end{center}

\vspace{1em}

% ============================================================
% Abstract
% ============================================================
\begin{abstract}
Gian-Carlo Rota proved that entropy is unique: every function on finite
probability distributions satisfying a short list of natural
properties---continuity, conditional additivity, maximality at
uniformity, and invariance under zero-probability outcomes---must be a
positive scalar multiple of Shannon entropy
$H(p_1,\dots,p_n) = -C\sum p_i \log p_i$.
This result, taught at MIT for decades but never formally published or
peer-reviewed, settles the question ``What is entropy?''\ that has
confounded physicists and information theorists since Clausius coined the
term in 1865.
We present the fully modernized proof of Rota's Entropy Theorem (RET),
machine-verified in Lean~4 with no \lean{sorry} and no custom axioms.
In the Lean formalization, Rota's original five properties become seven
explicit axioms, all of which are proven as theorems for the concrete
Shannon entropy function.
The capstone theorem,
\lean{RET\_All\_Entropy\_Is\_Scaled\_Shannon\_Entropy},
establishes that every entropy function satisfying these properties is
$C \times H_{\mathrm{Shannon}}$ for some $C > 0$.
Entropy is not disorder, not surprise, not tendency toward chaos.
Entropy is the \emph{exact count of information}---the number of bits
required to encode a system's state optimally---and Rota proved this is
the only possible such measure.
\end{abstract}

% ============================================================
\tableofcontents
\newpage

% ============================================================
% Section 1
% ============================================================
\section{Introduction: The Question Nobody Could Answer}
\label{sec:intro}

When Claude Shannon introduced his theory of communication in 1948, he
needed a name for his central quantity,
$H = -\sum_{i} p_i \log p_i$.
The legend, recounted by Myron Tribus, has John von Neumann advising
Shannon:

\begin{quote}
``You should call it entropy, for two reasons.  In the first place your
uncertainty function has been used in statistical mechanics under that
name, so it already has a name.  In the second place, and more
important, \emph{no one really knows what entropy is}, so in a debate
you will always have the advantage.''
\end{quote}

\noindent
Von Neumann's quip was not mere humor.  By 1948, ``entropy'' had
already accumulated a century of conflicting interpretations:

\begin{itemize}[nosep]
  \item \textbf{Thermodynamics} (Clausius, 1865): a quantity related to
    heat transfer and irreversibility, $dS = \delta Q / T$.
  \item \textbf{Statistical mechanics} (Boltzmann, 1877; Gibbs, 1902):
    a measure of the number of microstates consistent with a
    macrostate, $S = k_B \ln \Omega$.
  \item \textbf{Information theory} (Shannon, 1948):
    $H = -\sum p_i \log_2 p_i$, the minimum number of bits per symbol
    for lossless compression.
  \item \textbf{Quantum mechanics} (von Neumann, 1932):
    $S = -\mathrm{Tr}(\rho \ln \rho)$, the entropy of a density
    matrix.
\end{itemize}

Each community adopted its own notation, its own base of logarithm, its
own physical constant of proportionality, and its own informal
explanation.  Entropy has been variously described as ``disorder,''
``surprise,'' ``uncertainty,'' ``missing information,'' and ``tendency
toward chaos.''  None of these is a precise mathematical definition.
All of them are metaphors layered atop one another for 160 years.

The persistent confusion is not merely pedagogical.  It reflects a
genuine open question: \emph{Why is the same formula
$-\sum p_i \log p_i$ the right measure of information, of
thermodynamic irreversibility, and of quantum entanglement?}  Is this a
cosmic coincidence, or is there a mathematical reason?

Gian-Carlo Rota answered this question definitively.  He proved that the
formula is not a choice---it is a mathematical necessity.  Any function
satisfying a handful of self-evident properties \emph{must} be Shannon
entropy, up to a positive multiplicative constant.  All the entropies
listed above---Shannon's, Boltzmann's, Gibbs's, von Neumann's---are the
same function wearing different units.

Rota's proof was taught at MIT for decades, embedded in a
$\sim$400-page unpublished manuscript co-authored with Kenneth Baclawski
and Avner Billis.  It was never formally published.  It never went
through peer review.  And so it remained largely unknown to the broader
scientific community, even as physicists, computer scientists, and
engineers continued to debate what entropy ``really means.''

This paper aims to correct that record.

% ============================================================
% Section 2
% ============================================================
\section{The Man: Gian-Carlo Rota (1932--1999)}
\label{sec:rota}

Gian-Carlo Rota was born in Vigevano, Italy, in 1932.  His family fled
Mussolini's Italy in 1946, emigrating first to Ecuador and then to the
United States, where Rota enrolled at Princeton.  He completed his
Ph.D.\ at Yale in 1956 under the direction of Jacob T.~Schwartz, then
joined the faculty at MIT, where he remained for the rest of his career.

Rota is best known as the founder of modern algebraic combinatorics.
His 1964 paper ``On the Foundations of Combinatorial Theory~I: Theory of
M\"obius Functions'' transformed a collection of ad hoc counting tricks
into a rigorous algebraic theory, launching an entirely new field.  He
won the Steele Prize in 1988 and was elected to the National Academy of
Sciences in 1982.

Less well known, but equally significant, is Rota's place in the
intellectual lineage from von Neumann.  Von Neumann was best friends and, 
at the Manhattan Project, intellectually inseparable from
Stanis\l{}aw Ulam; Ulam, in turn, profoundly influenced Rota during
Rota's time at Los Alamos.  The chain von Neumann $\to$ Ulam $\to$ Rota
represents three generations of mathematicians for whom probability,
information, and computation were not separate disciplines but facets of
a single structure.

For over two decades, Rota taught \textbf{18.313}---``Introduction to
Probability''---at MIT.  The course notes, co-authored with Baclawski
and Billis, grew into a massive manuscript.  Chapter~7 of the 1979
edition (Chapter~8 of the 1992 revision) contains the entropy uniqueness
theorem.  Rota himself warned the reader:

\begin{quote}
``The proof is rather technical, so we suggest omitting it on the first
reading.'' \hfill ---Rota, Baclawski, Billis (\textit{Introduction to Probability
and Random Processes}, Ch.~8)
\end{quote}

\noindent
The manuscript was never published by a press, never submitted to a
journal, and never subjected to formal peer review.  Copies circulated
among Rota's students and colleagues, but the entropy theorem remained
hidden in plain sight---taught to generations of MIT students, yet
invisible to the broader community debating the meaning of entropy.

Rota died on April~18, 1999, at the age of 66.  His entropy theorem
deserves to be part of the scientific record.  This paper, and its
accompanying machine-verified Lean~4 formalization, aims to put it
there.

% ============================================================
% Section 3
% ============================================================
\section{Rota's Five Properties of Entropy}
\label{sec:properties}

Let $\Pi_n$ denote the set of \emph{probability distributions} on $n$
outcomes: vectors $(p_1, \dots, p_n)$ with $p_i \geq 0$ and
$\sum_{i=1}^{n} p_i = 1$.  Let $\Hfunc$ be a function that assigns a
non-negative real number to every finite probability distribution.
Rota identified five properties that any reasonable entropy function
should satisfy.

\begin{axiom_def}[Defined on probability distributions]
\label{ax:prob}
$\Hfunc$ is defined on $\Pi_n$ for every $n \geq 1$.  That is, $\Hfunc$
takes as input a finite probability distribution summing to~$1$ and
returns a non-negative real number.
\end{axiom_def}

\begin{axiom_def}[Zero invariance]
\label{ax:zero}
Adding outcomes with probability~$0$ does not change the entropy:
\[
  \Hfunc(p_1, \dots, p_n, 0) = \Hfunc(p_1, \dots, p_n).
\]
\end{axiom_def}

\begin{axiom_def}[Continuity]
\label{ax:cont}
$\Hfunc$ is a continuous function of its arguments.  That is, for any
distribution $\mathbf{p} \in \Pi_n$ and any $\varepsilon > 0$, there
exists $\delta > 0$ such that if
$|q_i - p_i| < \delta$ for all $i$, then
$|\Hfunc(\mathbf{q}) - \Hfunc(\mathbf{p})| < \varepsilon$.
\end{axiom_def}

\begin{axiom_def}[Conditional additivity]
\label{ax:add}
If a system is decomposed into a prior distribution $\sigma$ and
conditional distributions $\pi_i$ for each outcome $i$ of $\sigma$,
then:
\[
  \Hfunc(\text{joint}) = \Hfunc(\sigma) +
  \sum_{i} \sigma_i \cdot \Hfunc(\pi_i).
\]
This is the chain rule for entropy: the information in the whole equals
the information in the first stage plus the expected information in the
second.
\end{axiom_def}

\begin{axiom_def}[Maximum at uniform distribution]
\label{ax:max}
Among all distributions on $n$ outcomes, $\Hfunc$ is maximized by the
uniform distribution:
\[
  \Hfunc(p_1, \dots, p_n) \leq
  \Hfunc\!\left(\tfrac{1}{n}, \dots, \tfrac{1}{n}\right).
\]
\end{axiom_def}

\begin{remark}
Rota's Axioms~\ref{ax:prob}--\ref{ax:max} correspond to those in the
1979 edition (Chapter~7) and the 1992 revision (Chapter~8) of the
Rota--Baclawski manuscript.  The statement of conditional additivity
(Axiom~\ref{ax:add}) is sometimes written as
$H(\pi \mid \sigma) = H(\pi) - H(\sigma)$ in Rota's notation, which is
equivalent to the chain rule form above.  The key insight is that these
five properties \emph{uniquely determine} $\Hfunc$ up to a positive
multiplicative constant.
\end{remark}

% ============================================================
% Section 4
% ============================================================
\section{The Proof: From Axioms to Uniqueness}
\label{sec:proof}

Rota's proof proceeds in six stages.  We outline each, noting the
corresponding Lean theorem names in the formalization.

\subsection{Stage 1: Entropy of certainty is zero}

Define $f(n) = \Hfunc(1/n, \dots, 1/n)$---the entropy of the uniform
distribution on $n$ outcomes.  From Axioms~\ref{ax:prob} and
\ref{ax:zero}, we immediately obtain:
\[
  f(1) = 0.
\]
A distribution concentrated on a single outcome carries no information.

\medskip\noindent
\textbf{Lean:} \lean{f0\_1\_eq\_0}

\subsection{Stage 2: Monotonicity}

From Axiom~\ref{ax:max} (maximum at uniform) and Axiom~\ref{ax:zero}
(zero invariance), we deduce that $f(n)$ is monotone non-decreasing:
\[
  n \leq m \implies f(n) \leq f(m).
\]
More outcomes means at least as much entropy.

\medskip\noindent
\textbf{Lean:} \lean{f0\_mono}

\subsection{Stage 3: The Cauchy functional equation}

This is the crucial algebraic step.  Consider the product of two
uniform distributions: one on $m$ outcomes and one on $n$ outcomes.  The
joint distribution is uniform on $m \cdot n$ outcomes.  By the
conditional additivity axiom (Axiom~\ref{ax:add}):
\[
  f(m \cdot n) = f(m) + f(n).
\]
This is the Cauchy functional equation for the function $f$ restricted
to the positive integers.

\medskip\noindent
\textbf{Lean:} \lean{f0\_mul\_eq\_add\_f0}

\subsection{Stage 4: The logarithmic trapping argument}

This is the heart of Rota's proof and its most technically demanding
passage.  Fix $b \geq 2$ (we take $b = 2$) and consider any $n \geq 2$.
For any positive integer $m$, there exists a unique integer $k$ such
that:
\[
  b^k \leq n^m < b^{k+1}.
\]
Taking logarithms: $k \leq m \log_b n < k + 1$, so
\[
  \frac{k}{m} \leq \log_b n < \frac{k+1}{m}.
\]

Now apply the monotonicity and power law properties.  From $b^k \leq n^m$
and monotonicity of $f$:
\[
  k \cdot f(b) \leq m \cdot f(n).
\]
From $n^m < b^{k+1}$ and monotonicity:
\[
  m \cdot f(n) \leq (k+1) \cdot f(b).
\]

If $f(b) > 0$ (i.e., $\Hfunc$ is non-trivial), dividing through yields:
\[
  \frac{k}{m} \leq \frac{f(n)}{f(b)} \leq \frac{k+1}{m}.
\]

Both $f(n)/f(b)$ and $\log_b n$ are trapped in the same interval of
width $1/m$.  Since $m$ is arbitrary, letting $m \to \infty$ forces:
\[
  \frac{f(n)}{f(b)} = \log_b n
  \qquad\Longrightarrow\qquad
  f(n) = f(b) \cdot \log_b n = C \cdot \log n,
\]
where $C = f(b) / \log b$ is a positive constant independent of $n$.

\medskip\noindent
\textbf{Lean:} \lean{logarithmic\_trapping},
\lean{uniformEntropy\_power\_law}

\subsection{Stage 5: Extension to rational probabilities}

For a distribution with rational probabilities
$p_i = a_i / N$ where $\sum a_i = N$, the conditional additivity axiom
yields:
\[
  \Hfunc(p_1, \dots, p_n) = f(N) - \sum_{i} p_i \cdot f(a_i)
  = C \left(\log N - \sum_{i} p_i \log a_i\right)
  = C \sum_{i} p_i \log \frac{1}{p_i}.
\]
This is precisely $C$ times Shannon entropy.

\medskip\noindent
\textbf{Lean:} \lean{RUE\_rational\_case}

\subsection{Stage 6: Extension to real probabilities}

By the continuity axiom (Axiom~\ref{ax:cont}), the result for rational
probabilities extends to all real-valued probabilities by density of
$\mathbb{Q}$ in $\mathbb{R}$.  We conclude:

\begin{theorem}[Rota's Entropy Theorem---Uniqueness]
\label{thm:ret}
Let $\Hfunc$ be a function on finite probability distributions
satisfying Axioms~\ref{ax:prob}--\ref{ax:max}.  Then there exists a
constant $C \geq 0$ such that for all $(p_1, \dots, p_n) \in \Pi_n$:
\[
  \Hfunc(p_1, \dots, p_n) = C \sum_{i=1}^{n} p_i \log \frac{1}{p_i}.
\]
If $\Hfunc$ is non-trivial (i.e., not identically zero), then $C > 0$.
\end{theorem}

\medskip\noindent
\textbf{Lean:} \lean{RotaUniformTheorem},
\lean{RotaUniformTheorem\_formula\_with\_C\_constant}

% ============================================================
% Section 5
% ============================================================
\section{The Modernization: From Pen-and-Paper to Machine-Verified Proof}
\label{sec:lean}

The formalization was carried out in Lean~4 using the mathlib4 library.
It lives in the \lean{EGPT/Entropy/} module and spans three files:
\lean{Common.lean} (axiom definitions), \lean{RET.lean} (the abstract
uniqueness proof), and \lean{H.lean} (concrete verification that Shannon
entropy satisfies all axioms).  A fourth file,
\lean{NumberTheory/Analysis.lean}, provides the capstone theorem.

\subsection{From five properties to seven axioms}

Rota's five pen-and-paper properties translate into seven explicit Lean
axioms bundled in the structure \lean{HasRotaEntropyProperties}:

\begin{center}
\begin{tabular}{@{}rll@{}}
\toprule
\textbf{\#} & \textbf{Axiom (Lean name)} & \textbf{Origin} \\
\midrule
1 & \lean{IsEntropyNormalized} & Implicit in Rota (from Ax.~\ref{ax:prob}+\ref{ax:zero}) \\
2 & \lean{IsEntropySymmetric} & Implicit (Rota treats distributions as sets) \\
3 & \lean{IsEntropyZeroOnEmptyDomain} & Lean type-system edge case ($\mathrm{Fin}\;0$) \\
4 & \lean{IsEntropyZeroInvariance} & Axiom~\ref{ax:zero} \\
5 & \lean{IsEntropyContinuous} & Axiom~\ref{ax:cont} \\
6 & \lean{IsEntropyCondAddSigma} & Axiom~\ref{ax:add} \\
7 & \lean{IsEntropyMaxUniform} & Axiom~\ref{ax:max} \\
\bottomrule
\end{tabular}
\end{center}

The first two axioms were implicit in Rota's presentation.
\emph{Normalization} ($\Hfunc(\{1\}) = 0$) follows trivially on paper
from the zero-invariance and probability-set axioms, but making it
explicit simplifies the formal proof chain.  \emph{Symmetry}---invariance
under relabeling of outcomes---was automatic for Rota because he treated
distributions as unordered sets; in Lean, where $\mathrm{Fin}\;n \to
\NNR$ is an ordered function type, symmetry must be stated
explicitly as invariance under composition with an equivalence
$\alpha \simeq \beta$.  The third axiom, $\Hfunc(\emptyset) = 0$,
handles the edge case $\mathrm{Fin}\;0$ that Rota's framework never
considered.

\subsection{The abstract uniqueness proof (\texttt{RET.lean})}

File \lean{RET.lean} is a faithful rendition of Rota's proof.  It
assumes an abstract \lean{H\_func} satisfying
\lean{HasRotaEntropyProperties} and derives the key results listed in
Section~\ref{sec:proof}:

\begin{itemize}[nosep]
  \item \lean{f0\_1\_eq\_0}: $f(1) = 0$.
  \item \lean{f0\_mono}: $f$ is monotone non-decreasing.
  \item \lean{f0\_mul\_eq\_add\_f0}: $f(mn) = f(m) + f(n)$.
  \item \lean{uniformEntropy\_power\_law}: $f(n^k) = k \cdot f(n)$.
  \item \lean{logarithmic\_trapping}: the trapping argument forcing
    $f(n)/f(b) = \log_b n$.
  \item \lean{RotaUniformTheorem}: $f(n) = C \cdot \log n$.
\end{itemize}

The structure \lean{EntropyFunction} bundles an \lean{H\_func} together
with a proof that it satisfies \lean{HasRotaEntropyProperties}:

\begin{verbatim}
  structure EntropyFunction where
    (H_func : forall {a : Type} [Fintype a],
              (a -> NNReal) -> NNReal)
    (props : HasRotaEntropyProperties H_func)
\end{verbatim}

\subsection{Proving the axioms as theorems (\texttt{H.lean})}

The modernization step that goes beyond Rota's original work is
verifying that the concrete Shannon entropy function satisfies all seven
axioms.  File \lean{H.lean} defines the canonical entropy function
\lean{H\_canonical\_ln}---standard Shannon entropy using the natural
logarithm---and proves each axiom as a theorem:

\begin{enumerate}[nosep]
  \item \lean{h\_canonical\_is\_normalized}: $H(\{1\}) = 0$.
  \item \lean{h\_canonical\_is\_symmetric}: $H$ is invariant under
    relabeling.
  \item \lean{h\_canonical\_is\_zero\_on\_empty}: $H(\emptyset) = 0$.
  \item \lean{h\_canonical\_is\_zero\_invariance}: appending a
    zero-probability outcome does not change $H$.
  \item \lean{h\_canonical\_is\_continuous}: $H$ is continuous in
    $\varepsilon$-$\delta$ sense.
  \item \lean{h\_canonical\_is\_cond\_add\_sigma}: the chain rule
    holds.
  \item \lean{h\_canonical\_is\_max\_uniform}: $H$ is maximized by
    the uniform distribution.
\end{enumerate}

These are bundled into a single witness:

\begin{verbatim}
  noncomputable def TheCanonicalEntropyFunction_Ln :
    EntropyFunction := { ... }
\end{verbatim}

\noindent
This instance certifies that Shannon entropy is a valid
\lean{EntropyFunction}---and by the uniqueness theorem, the
\emph{only} one (up to scaling).

\subsection{The capstone theorem}

The culmination appears in \lean{NumberTheory/Analysis.lean}:

\begin{verbatim}
  theorem RET_All_Entropy_Is_Scaled_Shannon_Entropy
    (ef : EntropyFunction) (C : R) (hC_pos : 0 < C) :
    HasRotaEntropyProperties
      (fun p => (C * (ef.H_func p : R)).toNNReal)
\end{verbatim}

\noindent
This says: given \emph{any} entropy function \lean{ef} satisfying the
seven axioms, and \emph{any} positive constant $C$, the scaled function
$C \cdot H$ again satisfies all seven axioms.  Combined with the
uniqueness result from \lean{RET.lean}, this establishes that the space
of entropy functions is the one-dimensional ray
$\{C \cdot H_{\mathrm{Shannon}} : C > 0\}$.

The entire chain---from axiom definitions through the abstract
uniqueness proof to the concrete verification and the capstone---is
\lean{sorry}-free and uses no custom axioms.

As a final ``unit test,'' the formalization proves:

\begin{verbatim}
  theorem entropy_of_fair_coin_is_one_bit :
    H_canonical_log2 (canonicalUniformDist 2 ...) = 1
\end{verbatim}

\noindent
A fair coin flip has exactly one bit of entropy.  The engine is
calibrated.

% ============================================================
% Section 6
% ============================================================
\section{What Entropy Actually Is}
\label{sec:what}

With Rota's theorem in hand, we can state plainly what entropy is and
what it is not.

\paragraph{Entropy is not disorder.}
The ``disorder'' metaphor, ubiquitous in textbooks, has no mathematical
content.  It confuses a human judgment (``this looks messy'') with a
precisely defined quantity.  A shuffled deck of cards and a sorted deck
occupy the same state space; neither is more or less ``disordered'' in
any formal sense.

\paragraph{Entropy is not surprise.}
The ``surprise'' interpretation---common in information theory
pedagogy---captures intuition about rare events carrying more
information.  But surprise is a property of a single outcome;
entropy is a property of an entire distribution.

\paragraph{Entropy is not tendency toward chaos.}
This conflates a state function (entropy) with a dynamical law (the
second law of thermodynamics).  Entropy does not ``tend'' anywhere.
The second law says the entropy of an isolated system does not decrease;
entropy itself is just a number attached to a probability distribution.

\paragraph{Entropy is the exact count of information.}
RET proves that the \emph{unique} function satisfying
Axioms~\ref{ax:prob}--\ref{ax:max} is:
\[
  \Hfunc(p_1, \dots, p_n) = C \sum_{i=1}^{n} p_i \log \frac{1}{p_i},
\]
for some constant $C > 0$.  When $C = 1$ and the logarithm is base~2,
this is the number of bits required to encode a random variable
optimally.  When $C = k_B$ and the logarithm is natural, this is
Boltzmann's entropy.  The formula is the same; only the units change.

Entropy counts information.  The logarithm is not an arbitrary
choice---it is a mathematical \emph{necessity}, forced by the
monotonicity and multiplicative-to-additive property of $f$.  Any
monotone function satisfying $f(mn) = f(m) + f(n)$ is a logarithm; the
trapping argument of Section~\ref{sec:proof} makes this rigorous.

\paragraph{The double-entry accounting metaphor.}
If truth is the record of what actually happened in the past, then
entropy is its ledger.  Every physical process that resolves uncertainty
creates information; entropy is the quantity that ensures the books
balance.  Rota's theorem is the double-entry accounting system: it
guarantees that there is exactly one way to keep the ledger, up to a
choice of currency (the constant $C$).  You may measure in bits, nats,
or joules per kelvin, but the ledger is the same.

\paragraph{All entropies are one.}
Shannon entropy, Boltzmann entropy, Gibbs entropy, von Neumann entropy,
R\'enyi entropy of order~1---all are the same function
$C \sum p_i \log(1/p_i)$ with different values of $C$ and different
bases of logarithm.  RET unifies them.

% ============================================================
% Section 7
% ============================================================
\section{Implications}
\label{sec:implications}

\subsection{Einstein's question answered}

Einstein spent decades troubled by the tension between the continuous
fields of general relativity and the discrete quanta of quantum
mechanics.  If entropy is the unique measure of information, and
information is the same whether encoded in continuous or discrete form,
then the apparent tension dissolves.

The Lean formalization proves that the continuous Shannon entropy
$H_{\mathrm{continuous}}$ and the discrete Shannon entropy
$H_{\mathrm{discrete}}$ are related by
$H_{\mathrm{continuous}} = C \times H_{\mathrm{discrete}}$ for a
universal constant $C$.  Continuous and discrete descriptions are not
different physics---they are different encodings of the same information.

\subsection{The AI energy question}

If entropy equals computation, and computation is counting, then the
optimal encoding of computation is exact arithmetic---integer operations
(IOPs) rather than floating-point operations (FLOPs).  Floating-point
arithmetic introduces rounding error at every step; error accumulation
then demands larger word sizes, more memory, and more energy to
compensate.

The EGPT integer-only math library (EGPTMath) demonstrates this
principle: exact arithmetic using prime-power factorized number
representations replaces every FLOP with an IOP.  The Shannon-coding
interpretation of entropy guarantees that this encoding is optimal---no
bits are wasted, and no error accumulates.  The implication for AI
hardware is clear: if computation can be made exact, the energy cost of
error correction vanishes.

\subsection{Physics unification via statistical distributions}

The Lean formalization extends beyond the abstract uniqueness theorem.
The accompanying \lean{Physics/} module proves that all three canonical
statistical distributions---Bose--Einstein, Fermi--Dirac, and
Maxwell--Boltzmann---satisfy $H = C \times H_{\mathrm{Shannon}}$ over
$\mathbb{R}$.  This means that the entropy of \emph{every} physical
system governed by quantum or classical statistics is a scaled Shannon
entropy.  There is no separate ``physical'' entropy and ``information''
entropy; there is only entropy.

% ============================================================
% Section 8
% ============================================================
\section{Conclusion: Credit Where It Is Due}
\label{sec:conclusion}

Gian-Carlo Rota solved the entropy question decades ago.  His proof
established, with mathematical certainty, that entropy is unique: every
function satisfying a small number of natural axioms is a constant
multiple of Shannon entropy.  The logarithm is not a convention but a
theorem.  The proliferating definitions of entropy across
physics, information theory, and computer science are not competing
theories but different units of the same measure.

The result was available but hidden.  Buried in an unpublished
manuscript, taught to MIT students but never submitted for publication,
Rota's entropy theorem missed its moment.  The scientific community
continued debating what entropy means while the answer sat in a filing
cabinet in Cambridge, Massachusetts.

This paper and its accompanying Lean~4 formalization aim to correct the
record.  The proof is now machine-verified, permanent, and public.
Every axiom is a proven theorem.  Every step from axioms to uniqueness
has been checked by a computer.  There are no \lean{sorry}s, no custom
axioms, no gaps.

The intellectual debt runs deep.  Von Neumann, who told Shannon that
``no one really knows what entropy is,'' was one mind at the Manhattan project with Ulam, who influenced
Rota, who proved that the answer was always implicit in the question.
Entropy is the count of information.  It always was.

It is time the world knew.

% ============================================================
% Acknowledgments
% ============================================================
\section*{Acknowledgments}

The author thanks the Lean community and the mathlib4 contributors
whose library made the formalization possible.  The original ideas
belong to Gian-Carlo Rota, Kenneth Baclawski, and Avner Billis.  This
paper merely formalizes and publicizes what they created.

% ============================================================
% References
% ============================================================
\begin{thebibliography}{99}

\bibitem{rota-baclawski}
G.-C.~Rota, K.~Baclawski, and A.~Billis,
\textit{Introduction to Probability and Random Processes},
unpublished manuscript, MIT, 1979 (revised 1992/1993).
Chapter~7 (1979 edition) / Chapter~8 (1992 edition):
``The Uniqueness of Entropy.''

\bibitem{shannon1948}
C.~E.~Shannon,
``A Mathematical Theory of Communication,''
\textit{Bell System Technical Journal}, vol.~27, pp.~379--423, 623--656,
1948.

\bibitem{khinchin1957}
A.~I.~Khinchin,
\textit{Mathematical Foundations of Information Theory},
Dover Publications, New York, 1957.
(Translation of two Russian papers from 1953 and 1956.)

\bibitem{vonneumann1932}
J.~von Neumann,
\textit{Mathematische Grundlagen der Quantenmechanik},
Springer, Berlin, 1932.
(English translation: \textit{Mathematical Foundations of Quantum
Mechanics}, Princeton University Press, 1955.)

\bibitem{vonneumann1958}
J.~von Neumann,
\textit{The Computer and the Brain},
Yale University Press, New Haven, 1958.

\bibitem{rota1997}
G.-C.~Rota,
\textit{Indiscrete Thoughts},
Birkh\"auser, Boston, 1997.

\bibitem{boltzmann1877}
L.~Boltzmann,
``\"Uber die Beziehung zwischen dem zweiten Hauptsatze der mechanischen
W\"armetheorie und der Wahrscheinlichkeitsrechnung respektive den
S\"atzen \"uber das W\"armegleichgewicht,''
\textit{Wiener Berichte}, vol.~76, pp.~373--435, 1877.

\bibitem{gibbs1902}
J.~W.~Gibbs,
\textit{Elementary Principles in Statistical Mechanics},
Charles Scribner's Sons, New York, 1902.

\bibitem{clausius1865}
R.~Clausius,
``Ueber verschiedene f\"ur die Anwendung bequeme Formen der
Hauptgleichungen der mechanischen W\"armetheorie,''
\textit{Annalen der Physik}, vol.~201, no.~7, pp.~353--400, 1865.

\bibitem{lean4formalization}
E.~Abadir,
\textit{EGPT: Electronic Graph Paper Theory---Lean 4 Formalization},
\url{https://eabadir.github.io/EGPT},
2024--2026.
Files: \texttt{Lean/EGPT/Entropy/\{Common,RET,H\}.lean},
\texttt{Lean/EGPT/NumberTheory/Analysis.lean}.

\bibitem{rota1964}
G.-C.~Rota,
``On the Foundations of Combinatorial Theory~I: Theory of M\"obius
Functions,''
\textit{Zeitschrift f\"ur Wahrscheinlichkeitstheorie und Verwandte
Gebiete}, vol.~2, pp.~340--368, 1964.

\bibitem{tribus1971}
M.~Tribus and E.~C.~McIrvine,
``Energy and Information,''
\textit{Scientific American}, vol.~225, no.~3, pp.~179--188, 1971.
(Source of the von Neumann--Shannon anecdote.)

\end{thebibliography}

% ============================================================
% Appendix
% ============================================================
\appendix

\section{Summary of Lean Theorem Names}
\label{app:lean}

For reference, the complete chain of formal results:

\bigskip

\begin{center}
\begin{tabular}{@{}ll@{}}
\toprule
\textbf{Lean identifier} & \textbf{Statement} \\
\midrule
\multicolumn{2}{@{}l}{\textit{Abstract uniqueness (\texttt{RET.lean})}} \\
\lean{f0\_1\_eq\_0} & $f(1) = 0$ \\
\lean{f0\_mono} & $f$ is monotone \\
\lean{f0\_mul\_eq\_add\_f0} & $f(mn) = f(m) + f(n)$ \\
\lean{uniformEntropy\_power\_law} & $f(n^k) = k \cdot f(n)$ \\
\lean{logarithmic\_trapping} & Trapping argument \\
\lean{RotaUniformTheorem} & $\exists C,\; f(n) = C \log n$ \\[3pt]
\midrule
\multicolumn{2}{@{}l}{\textit{Concrete axiom proofs (\texttt{H.lean})}} \\
\lean{h\_canonical\_is\_normalized} & $H(\{1\}) = 0$ \\
\lean{h\_canonical\_is\_symmetric} & Relabeling invariance \\
\lean{h\_canonical\_is\_zero\_on\_empty} & $H(\emptyset) = 0$ \\
\lean{h\_canonical\_is\_zero\_invariance} & Zero-probability invariance \\
\lean{h\_canonical\_is\_continuous} & $\varepsilon$-$\delta$ continuity \\
\lean{h\_canonical\_is\_cond\_add\_sigma} & Chain rule \\
\lean{h\_canonical\_is\_max\_uniform} & Maximum at uniform \\[3pt]
\midrule
\multicolumn{2}{@{}l}{\textit{Bundled witness (\texttt{H.lean})}} \\
\lean{TheCanonicalEntropyFunction\_Ln} & Shannon entropy is an \lean{EntropyFunction} \\[3pt]
\midrule
\multicolumn{2}{@{}l}{\textit{Capstone (\texttt{Analysis.lean})}} \\
\lean{RET\_All\_Entropy\_Is\_Scaled\_Shannon\_Entropy} & $\forall C>0$: $C \cdot H$ satisfies all axioms \\[3pt]
\midrule
\multicolumn{2}{@{}l}{\textit{Calibration (\texttt{H.lean} / \texttt{Analysis.lean})}} \\
\lean{entropy\_of\_fair\_coin\_is\_one\_bit} & $H_{\log_2}(\text{fair coin}) = 1$ bit \\
\bottomrule
\end{tabular}
\end{center}

\end{document}
